\SourcePage{311}{316}
\section{The unitarity condition}

The scattering matrix must be unitary: $\hat S\hat S^+=1$, or in matrix elements,
\begin{equation}
(SS^+)_{fi}=\sum_n S_{fn}S_{in}^*=\delta_{fi},
\tag{71.1}
\end{equation}
where $n$ labels all possible intermediate states.\footnote{The meaning of $\delta_{fi}$ in (71.1) depends, of course, on the particular choice of quantum numbers and on the normalization of the system's wave functions. It must be defined so that $\sum_f\delta_{if}=1$.} This is the most general property of the $S$-matrix and

\SourcePage{312}{317}
ensures preservation of the normalization and orthogonality of states in a reaction (cf. vol.~III, \S\S~125, 144). In particular, the diagonal elements of (71.1) state simply that the sum of transition probabilities from a given initial state to every possible final state is unity:
\[
\sum_n\lvert S_{ni}\rvert^2=1.
\]

Substitution of (64.2) into (71.1) gives
\begin{equation}
\begin{aligned}
T_{fi}-T_{if}^*&=i(2\pi)^4\sum_n\delta^{(4)}(P_f-P_n)T_{fn}T_{in}^*={}\\
&=i(2\pi)^4\sum_n\delta^{(4)}(P_f-P_n)T_{nf}^*T_{ni}.
\end{aligned}
\tag{71.2}
\end{equation}
The two equivalent forms on the right follow by writing unitarity as $\hat S\hat S^+=1$ or $\hat S^+\hat S=1$, with the factors in opposite orders.

The left side is linear, whereas the right side is quadratic, in the matrix elements $T$. Thus, if the interaction contains a small parameter, as does the electromagnetic interaction, the two sides are respectively of first and second order. In the first approximation the right side may be neglected, giving
\begin{equation}
T_{fi}=T_{if}^*,
\tag{71.3}
\end{equation}
so the matrix $T$ is Hermitian.

To make (71.2) more explicit, the sum over $n$ must be specified. Consider a two-particle collision for which conservation laws permit only elastic scattering; all intermediate states in (71.2) are then two-particle states. Summation over them means integration over the intermediate momenta $\mathbf p_1'',\mathbf p_2''$ and summation over the spin quantum numbers, such as helicities, of both particles, denoted collectively by $\lambda''$:
\[
\sum_n\longrightarrow\int\frac{V^2d^3p_1''d^3p_2''}{(2\pi)^6}\sum_{\lambda''}.
\]
Eliminating the delta functions as in \S~64 gives the two-particle unitarity condition
\[
T_{fi}-T_{if}^*=\frac{iV^2}{(2\pi)^2}\sum_{\lambda''}\frac{\lvert\mathbf p\rvert}{\varepsilon}
\int T_{fn}T_{in}^*\varepsilon_1''\varepsilon_2''\,do'',
\]

\SourcePage{313}{318}
where $\mathbf p$ and $\varepsilon$ are the momentum and total energy in the center-of-momentum frame. On passing from $T_{fi}$ to $M_{fi}$ by (64.10), the normalization volume disappears:
\begin{equation}
M_{fi}-M_{if}^*=\frac{i}{(4\pi)^2}\sum_{\lambda''}\frac{\lvert\mathbf p\rvert}{\varepsilon}
\int M_{fn}M_{in}^*\,do''.
\tag{71.4}
\end{equation}

Define the elastic-scattering amplitude by
\begin{equation}
d\sigma=\lvert\langle\mathbf n'\lambda'\vert f\vert\mathbf n\lambda\rangle\rvert^2do'
\tag{71.5}
\end{equation}
($\mathbf n,\mathbf n'$ are the initial and final momentum directions, and $\lambda,\lambda'$ the initial and final spin quantum numbers). Comparison with (64.19) gives
\begin{equation}
\langle\mathbf n'\lambda'\vert f\vert\mathbf n\lambda\rangle=\frac{1}{8\pi\varepsilon}M_{fi},
\tag{71.6}
\end{equation}
and (71.4) becomes
\begin{equation}
\begin{aligned}
&\langle\mathbf n'\lambda'\vert f\vert\mathbf n\lambda\rangle
-\langle\mathbf n\lambda\vert f\vert\mathbf n'\lambda'\rangle^*={}\\
&\qquad=\frac{i\lvert\mathbf p\rvert}{2\pi}\sum_{\lambda''}\int
\langle\mathbf n'\lambda'\vert f\vert\mathbf n''\lambda''\rangle
\langle\mathbf n\lambda\vert f\vert\mathbf n''\lambda''\rangle^*\,do'',
\end{aligned}
\tag{71.7}
\end{equation}
which generalizes the familiar nonrelativistic formula (see vol.~III, (125.8)).

The forward elastic-scattering amplitude is the diagonal matrix element $T_{ii}$ for which the final particle state coincides with the initial state.\footnote{This means an element of $T$, not of $S$: the diagonal element is taken after the unit matrix has been removed from $S$.} For this amplitude, (71.2) gives
\begin{equation}
2\Im T_{ii}=(2\pi)^4\sum_n\lvert T_{in}\rvert^2\delta^{(4)}(P_i-P_n).
\tag{71.8}
\end{equation}
The right side differs only by a factor from the total cross section $\sigma_t$ of all possible scattering processes from the initial state $i$. Indeed, summing (64.5) over $f$ and dividing by the flux density $j$ gives
\[
\sigma_t=\frac{(2\pi)^4V}{j}\sum_n\lvert T_{in}\rvert^2\delta^{(4)}(P_i-P_n),
\]
so that
\[
\frac{2V}{j}\Im T_{ii}=\sigma_t.
\]
Replacing $T_{ii}=M_{ii}/(2\varepsilon_1V\cdot2\varepsilon_2V)$ and using $j$ from (64.17) removes the normalization volume:
\begin{equation}
\Im M_{ii}=2\lvert\mathbf p\rvert\varepsilon\sigma_t.
\tag{71.9}
\end{equation}

\SourcePage{314}{319}
This is the content of the \emph{optical theorem}. In terms of the elastic amplitude (71.6), it takes its usual form
\begin{equation}
\Im\langle\mathbf n\lambda\vert f\vert\mathbf n\lambda\rangle
=\frac{\lvert\mathbf p\rvert}{4\pi}\sigma_t
\tag{71.10}
\end{equation}
(cf. vol.~III, (142.10)).

If the $S$-matrix is given in the angular-momentum representation, as partial amplitudes, its diagonality in $J$ allows the unitarity condition to be written separately for each $J$.

If only elastic scattering is possible,
\begin{equation}
\sum_{\lambda''}\langle\lambda'\vert S^J\vert\lambda''\rangle
\langle\lambda\vert S^J\vert\lambda''\rangle^*=\delta_{\lambda\lambda'}.
\tag{71.11}
\end{equation}
By $T$-invariance, the elastic-scattering matrix is symmetric (cf. (69.10)) and can therefore be diagonalized. Unitarity then requires its diagonal elements to have unit modulus; they are conventionally written
\begin{equation}
S_n^J=\exp(2i\delta_{Jn}),
\tag{71.12}
\end{equation}
where the real quantities $\delta_{Jn}$ are functions of the energy, and $n$ labels the diagonal elements for fixed $J$. In general, if the number $N$ of independent amplitudes exceeds the rank of the square matrix $S^J$, the coefficients of the transformation diagonalizing $S^J$ depend on $J,E$; together with its principal values they contain independent quantities equivalent to the original $N$. If $N$ equals the rank, and hence the number of principal values, the diagonalization coefficients are universal. The diagonalizing states are then states of definite parity, though no longer of definite helicity.

In terms of the partial amplitudes $\langle\lambda'\vert f^J\vert\lambda\rangle$, (71.11) is
\begin{equation}
\begin{aligned}
&\langle\lambda'\vert f^J\vert\lambda\rangle
-\langle\lambda\vert f^J\vert\lambda'\rangle^*={}\\
&\qquad=2i\lvert\mathbf p\rvert\sum_{\lambda''}
\langle\lambda'\vert f^J\vert\lambda''\rangle
\langle\lambda\vert f^J\vert\lambda''\rangle^*,
\end{aligned}
\tag{71.13}
\end{equation}
as follows by substituting (68.13) in (71.7) and using orthonormality of the $D$ functions. Under $T$-invariance the matrix $\langle\lambda'\vert f^J\vert\lambda\rangle$ is symmetric, and (71.13) becomes
\begin{equation}
\Im\langle\lambda'\vert f^J\vert\lambda\rangle
=\lvert\mathbf p\rvert\langle\lambda'\vert f^Jf^{J+}\vert\lambda\rangle.
\tag{71.14}
\end{equation}
If the matrix is diagonalized, its diagonal elements are
\begin{equation}
f_n^J=\frac{1}{2i\lvert\mathbf p\rvert}\bigl(\exp(2i\delta_{Jn})-1\bigr)
=\frac{1}{\lvert\mathbf p\rvert}\exp(i\delta_{Jn})\sin\delta_{Jn}.
\tag{71.15}
\end{equation}

\SourcePage{315}{320}
Finally, consider some consequences of unitarity together with $CPT$-invariance. The latter gives
\begin{equation}
T_{fi}=T_{\bar i\bar f},
\tag{71.16}
\end{equation}
where $\bar i,\bar f$ differ from $i,f$ by replacing every particle by its antiparticle and reversing angular-momentum vectors while leaving momenta unchanged. In particular,
\[
T_{ii}=T_{\bar i\bar i}.
\]
It follows from (71.8) or (71.9) that the total cross sections for all possible processes from specified initial states are the same for reactions of particles and antiparticles.

In particular, the total decay probabilities, and hence lifetimes, are the same for a particle and its antiparticle. Together with equality of their masses (\S~11), these results are among the most important consequences of $CPT$-invariance. Recall from the end of \S~69 that the same assertion separately for each decay channel also requires $CP$-invariance.

\ProblemHeading

Starting from unitarity, find the relation between the phases of the partial amplitudes for pion photoproduction on nucleons ($\gamma+N\to\pi+N$) and elastic pion--nucleon scattering ($\pi+N\to\pi+N$). Take into account that $\pi N$ scattering is due to the strong interaction, whereas photoproduction and $\gamma N$ scattering are electromagnetic.

\SolutionHeading Denote the partial amplitudes by
\[
\langle\pi N\vert S\vert\gamma N\rangle=S_{\pi\gamma},\qquad
\langle\gamma N\vert S\vert\gamma N\rangle=S_{\gamma\gamma},\qquad
\langle\pi N\vert S\vert\pi N\rangle=S_{\pi\pi}
\]
(the $J$ and helicity indices are suppressed). Photoproduction is first order and $\gamma N$ scattering second order in the charge $e$; hence $S_{\pi\gamma}\sim e$, $S_{\gamma\gamma}-1\sim e^2$. The amplitude $S_{\pi\pi}$ contains no small parameter. To terms of order $e$, (71.1) gives
\begin{align}
S_{\pi\gamma}S_{\gamma\gamma}^*+S_{\pi\pi}S_{\gamma\pi}^*
&\approx S_{\pi\gamma}+S_{\pi\pi}S_{\gamma\pi}^*=0,
\tag{1}\\
S_{\pi\gamma}S_{\pi\gamma}^*+S_{\pi\pi}S_{\pi\pi}^*
&\approx S_{\pi\pi}S_{\pi\pi}^*=1
\tag{2}
\end{align}
(the 1 on the right of (2) is the unit matrix in spin variables). By $T$-invariance, $S_{\pi\pi}$ is symmetric and $S_{\gamma\pi}=S_{\pi\gamma}$. Choose $S_{\pi\pi}$ diagonal, with respect to pion states of definite parity. Equation (2) then requires its diagonal elements to have the form $e^{2i\delta_\pi}$, with different constants $\delta_\pi$. Equation (1) gives for each element of $S_{\pi\gamma}$
\[
S_{\pi\gamma}/S_{\pi\gamma}^*=-e^{2i\delta_\pi},
\]
and hence
\[
S_{\pi\gamma}=\pm\lvert S_{\pi\gamma}\rvert i e^{i\delta_\pi}.
\]
Thus the phase of the partial photoproduction amplitude into a state of definite parity is determined by the phase of elastic $\pi N$ scattering.
